% !TeX root = ../main.tex
% Joint Solver (input into \section{Joint Solver} in main.tex).

Problem (P0) in \eqref{eq:P0} couples binary placement, integer port capacity,
continuous trajectories, a combinatorial assignment, and the nonconvex
finite-source wait \eqref{eq:wq}. Rather than attack this MINLP monolithically,
we decompose it into blocks that are each tractable and solve them in a loop:
a trajectory block, a finite-source queue evaluation, a port-allocation block,
an assignment block, and an outer placement search. Fig.~\ref{fig:solver}
sketches the flow. The blocks share the peak-AoI objective \eqref{eq:peak_m},
so a change accepted by one block is scored by the same criterion the others
use: the assignment search drives the peak AoI down monotonically to a local
optimum, while the placement and port searches are enumerated exhaustively and
retain the best configuration found. We describe each block, then the overall
algorithm.

\subsection{Trajectory block (GPU CETSP)}
\label{subsec:traj_block}
A UAV collects from any sensor within a horizontal radius $r_c$, so a patrol
need not visit every sensor point: it need only pass within $r_c$ of each
assigned sensor. The shortest patrol is therefore the shortest closed tour of
hover points whose $r_c$-disks cover the UAV's sub-field, a Close-Enough
Traveling Salesman Problem (CETSP)~\cite{cetsp}. We solve it by differentiable
optimization. For UAV $m$, let $\mathbf{p}_{m,1},\dots,\mathbf{p}_{m,n_m}$ be
its hover points in a fixed cyclic order (initialized by an angular sort of
$k$-means centroids of the sub-field, which keeps the closed-tour length
differentiable in the coordinates). Only the coordinates slide; the visiting
order is frozen. The objective is the total closed-tour length plus a coverage
penalty,
\begin{equation}
\begin{split}
  \mathcal{L}
  = \sum_{m\in\mathcal{M}}
    \Bigg[
    &\underbrace{\sum_{j=1}^{n_m}\big\|\mathbf{p}_{m,j{+}1}-\mathbf{p}_{m,j}\big\|}_{\text{tour length}}\\
    &+ \lambda_{\mathrm{cov}}
      \underbrace{\sum_{k\in\mathcal{K}_m}\big[\,d_{m,k}^{\min}-r_c\,\big]_+^2}_{\text{coverage violation}}\Bigg],
\end{split}
  \label{eq:cetsp_loss}
\end{equation}
where indices are cyclic ($\mathbf{p}_{m,n_m{+}1}\equiv\mathbf{p}_{m,1}$),
$\mathcal{K}_m$ is the sensor set assigned to UAV $m$, $d_{m,k}^{\min}$ is the
distance from sensor $k$ to its nearest hover point, $[\cdot]_+=\max(0,\cdot)$,
and $\lambda_{\mathrm{cov}}$ weights coverage against length. We minimize
\eqref{eq:cetsp_loss} by Adam over the hover-point coordinates. All UAVs of a
scenario are stacked into one padded, masked batch and optimized in parallel on
a GPU, so the whole swarm (and, by stacking seeds, whole sweeps) is optimized in
a single pass. The block returns each UAV's optimized tour length, hence its
flight time $t_{\mathrm{fly}}^{\mathrm{loop}}$ and, through \eqref{eq:taufly},
its revisit period $t_{\mathrm{loop}}$ and flight budget $\tau_{\mathrm{fly}}$.
Because the sub-fields are fixed by the field partition, this block is
independent of placement and is solved once at the start.

\subsection{Charging queue and port allocation}
\label{subsec:port_block}
Given an assignment, each open station $s$ is the finite-source queue of
Section~\ref{sec:model}. Its mean wait $\Wq(N_s,c_s)$ follows the machine-repair
balance equations \eqref{eq:pn} to \eqref{eq:wq} from the number of assigned UAVs $N_s$,
the ports $c_s$, the operating time \eqref{eq:uptime}, and the charge time
$\tau_{\mathrm{charge}}$; this evaluation is what turns a candidate configuration
into a peak AoI through \eqref{eq:peak_m}. Ports are then allocated across the
open stations to reduce the peak wait under the budget
$\sum_s c_s\le C_{\mathrm{tot}}$. Since $\Wq$ is convex-decreasing in the port
count, the peak-reducing allocation adds ports where they buy the most, that is,
to the most contended stations: a water-filling allocation. We realize it by
scanning the integer port splits that sum to $C_{\mathrm{tot}}$ and keeping the
split that yields the lowest peak AoI, which is exact for the small station
budgets $S_{\max}$ considered here.

\subsection{Assignment}
\label{subsec:assign_block}
With placement and ports fixed, UAVs are assigned to open stations to balance the
queueing load. Starting from the nearest-station assignment, we run a greedy
load-balancing local search: scanning the UAVs in turn, we apply the first
single-UAV reassignment that strictly reduces the peak AoI \eqref{eq:peak_m},
restart the scan after each accepted move, and stop when no single reassignment
improves it. Because the peak AoI is set by the worst
station, driving it down repeatedly pulls UAVs off the most contended station
onto lighter ones, so the search approximates the min-max (load-balancing)
assignment $\argmin_{a}\max_s \Wq(N_s,c_s)$ without enumerating all
assignments.

\subsection{Placement (outer search)}
\label{subsec:place_block}
The open-site subset $\mathbf{x}$ is chosen by an outer search over the candidate
combinations. Candidate sites are the interior nodes of a regular grid over the
field, and the search enumerates the $\binom{S}{S_{\max}}$ subsets of size
$S_{\max}$. For each subset, the inner blocks (port allocation, assignment, and
the queue evaluation) are run and the resulting peak AoI recorded; the subset
with the lowest peak AoI is kept. This outer search is where the placement flips
from coverage-driven to traffic-driven: a subset that spreads stations to split
the charging load can beat the minimum-detour subset once contention is priced
in through $\Wq$.

\subsection{Overall algorithm}
\label{subsec:algorithm}
Algorithm~\ref{alg:solver} assembles the blocks. The trajectory block is solved
once, after which the placement, port, and assignment blocks are alternated: for
each candidate placement the port split and the load-balancing assignment are
optimized against the same peak-AoI objective, each evaluated through the
finite-source queue.

\begin{algorithm}[t]
\caption{Block-coordinate joint solver for (P0)}
\label{alg:solver}
\begin{algorithmic}[1]
\REQUIRE sensors $\mathcal{K}$, candidate sites $\mathcal{S}$, station budget
  $S_{\max}$, port budget $C_{\mathrm{tot}}$, swarm size $M$
\ENSURE placement $\mathbf{x}$, ports $\mathbf{c}$, assignment $\mathbf{a}$,
  trajectories $\{\mathbf{q}_m\}$, peak AoI $\bar{A}$
\STATE partition $\mathcal{K}$ into $M$ sub-fields
\STATE \textbf{trajectory block:} minimize \eqref{eq:cetsp_loss} by Adam on the
  GPU for all UAVs; obtain $t_{\mathrm{loop},m}$, $\tau_{\mathrm{fly},m}$
\STATE $\bar{A}^\star \gets +\infty$
\FOR{each site subset $\mathbf{x}$ of size $S_{\max}$ in $\mathcal{S}$}
  \FOR{each port split $\mathbf{c}$ with $\sum_s c_s = C_{\mathrm{tot}}$}
    \STATE $\mathbf{a}\gets$ nearest-station assignment
    \REPEAT
      \STATE scan single-UAV reassignments; apply the first one that strictly
        reduces the peak AoI \eqref{eq:peak_m} (evaluating $\Wq$ by \eqref{eq:wq})
        and restart the scan
    \UNTIL{no single reassignment reduces the peak AoI}
    \STATE $\bar{A}\gets$ peak AoI of $(\mathbf{x},\mathbf{c},\mathbf{a})$
    \IF{$\bar{A} < \bar{A}^\star$}
      \STATE $\bar{A}^\star\gets\bar{A}$; store $(\mathbf{x},\mathbf{c},\mathbf{a})$
    \ENDIF
  \ENDFOR
\ENDFOR
\RETURN stored $(\mathbf{x},\mathbf{c},\mathbf{a})$, $\{\mathbf{q}_m\}$, $\bar{A}^\star$
\end{algorithmic}
\end{algorithm}

\emph{Complexity and convergence.} The trajectory block costs one batched
GPU optimization independent of the placement search. The outer loops evaluate
$\binom{S}{S_{\max}}$ placements times the number of port splits, and each
evaluation runs the greedy assignment, whose local search terminates because
every accepted move strictly lowers the peak AoI (a quantity bounded below by
zero) and there are finitely many assignments. For a fixed placement and port
split, the returned assignment is a local optimum of \eqref{eq:peak_m} under
single-UAV moves; because the placement and port searches are exhaustive over the
candidate grid and the integer splits, and the best configuration is retained,
the solver returns the lowest peak AoI it evaluates. We do not claim global
optimality: the trajectory order is fixed and the assignment search is greedy, so
the result is the best configuration reachable by this search from the
nearest-station initialization, not necessarily the global optimum of (P0).
